\documentclass[10pt,a4paper]{amsart}
\usepackage[margin=19mm]{geometry}
\usepackage{amsmath,amssymb,mathtools}
\usepackage[scr=rsfs,cal=euler]{mathalfa}
\usepackage{xfrac}
\usepackage[unicode,hidelinks,bookmarks=false]{hyperref}
\newcommand{\ie}{i.e.\ }
\newcommand{\cf}{cf.\ }
\newcommand{\resp}{respectively\ }
\newcommand{\Subsection}{\subsection}
\providecommand{\nobreakdash}{\nobreak\hbox{-}}
\setlength{\emergencystretch}{2em}
\multlinegap0pt
\usepackage{bm}
\usepackage{bbm}
\usepackage{dsfont}
\usepackage{xspace}
\usepackage{cancel}
\usepackage{mathtools}
\usepackage{amsthm}
\makeatletter
\@ifundefined{theorem}{\newtheorem{theorem}{Theorem}}{}
\@ifundefined{claim}{\newtheorem{claim}[theorem]{Claim}}{}
\@ifundefined{lemma}{\newtheorem{lemma}[theorem]{Lemma}}{}
\makeatother
\makeatletter
\renewcommand{\geq}{\geqslant}
\renewcommand{\leq}{\leqslant}
\expandafter\ifx\csname scproof\endcsname\relax
\newenvironment{scproof}[1][]{\begin{proof}[\textsc{\upshape{Proof}}#1]}{\end{proof}}
\fi
\makeatother
\title[Measurable one-ended spanning trees]{Measurable one-ended spanning trees}
\author{Matt Bowen, Ant\'onio Gir\~ao, H\'ector Jard\'on-S\'anchez, Grigory Terlov}
\date{Source: arXiv:2509.15000v2 (CC BY 4.0)}
\begin{document}
\maketitle

\noindent\emph{Source and licence.} Measurable one-ended spanning trees. Version arXiv:2509.15000v2 (CC BY 4.0), distributed under the Creative Commons Attribution 4.0 International licence, as recorded in the arXiv licence metadata for this exact version and shown on the abstract page https://arxiv.org/abs/2509.15000v2. Full rights notice: licenses/arxiv-2509.15000-CC-BY-4.0.txt.\par\smallskip

\noindent\textit{Editorial scope.} Counted unit: the Lemma whose source label is \texttt{lem:technical} in the article (source0.tex lines 438--444, source label \texttt{lem:technical}), delivered together with the article's own complete original proof of it (source0.tex lines 446--500, one \texttt{proof} environment of 55 lines). No mathematics is written, repaired or re-derived: statement and proof are the article's own text line for line. Required context retained uncounted: none. Every definition, display and label of those lines is retained unchanged. Omitted from this package and not redistributed: 1--437, 445--445, 501--775. Nothing inside a retained excerpt is omitted or altered: every retained body is delivered exactly as the article's lines are written, so no omission receipt of any kind accompanies this package. Citations: the retained excerpts carry no citation command at all, so no bibliography entry of the article is transcribed and no reference is redistributed. Reference adaptation: none was needed and none was made; every reference inside the delivered unit resolves inside it, so no unresolved reference or citation is produced. Printed slips kept literal: no printed word, symbol, hypothesis or number of the counted statement or of its proof was altered. Layout and class: the article's own class is \texttt{amsart} with packages including geometry, amssymb, graphicx, amsmath, tikz, hyperref, dsfont, lmodern, mathabx, comment, microtype, extdash, titlesec; the wrapper is the lane's pinned \texttt{amsart} editorial wrapper at 10pt on A4, which restates the theorem-like environments the retained text needs and adds the article's own macro definitions that the retained text needs (\texttt{\textbackslash geq}, \texttt{\textbackslash leq}), with extra wrapper packages (bm, bbm, dsfont, xspace, cancel, mathtools). The delivered numbering is the unit's own. Assets: no retained span contains a figure environment, a graphics-inclusion command, a PDF-page inclusion, a table or a TikZ picture; no image file is copied into this package, the package declares no asset dependency and no image file is redistributed with it. Licence: version arXiv:2509.15000v2 of this article is distributed under the Creative Commons Attribution 4.0 International licence, as recorded in the arXiv licence metadata for this exact version and shown on the abstract page https://arxiv.org/abs/2509.15000v2. A full rights notice is delivered at licenses/arxiv-2509.15000-CC-BY-4.0.txt. Characters: every retained excerpt is pure ASCII, so the delivered engine, XeLaTeX with its default Latin Modern text font, reports no missing character.
\begin{lemma}\label{lem:technical}
    Let $\mathcal{G}$ be a locally finite, one-ended, hyperfinite, mcp graph on a finite measure space $(X,\nu)$ and $\mathcal{G}' \subseteq \mathcal{G}$ a substantial measurable subgraph. Then, for any fber $\mathcal{F} \subseteq \mathcal{R}|_{\mathcal G[I(\mathcal G')]}$ and $\delta > 0$ there exists an acyclic Borel subgraph $\mathcal{T} \subseteq \mathcal{G}$ with finite connected components such that 
    \begin{itemize}
        \item[i)] $|[x]_{\mathcal{T}} \cap O(\mathcal{G}')| = 1$ for $\nu$-a.e.\ $x\in I(\mathcal{G}')$, and
        \item[ii)] $[x]_\mathcal{F} \subseteq [x]_{\mathcal{T}}$ for at least $(1 - \delta) \nu (I(\mathcal{G}'))$-many $x\in X$.
    \end{itemize}
\end{lemma}

\begin{proof}
    By hyperfiniteness of $\mathcal{G}$ there exists a fber $\mathcal E$ such that $\mathcal{F} \subset \mathcal{E} \subset \mathcal{R}|_{\mathcal G[I(\mathcal G')]}$, has $\mathcal G$-connected equivalence classes, each of whose intersection with $I(\mathcal G')$ is also connected, and  \begin{equation*}
    \nu \left(x \in I(\mathcal G') : \exists y \in O(\mathcal{G}')\ \text{such that}\ (y,x) \in \mathcal{E}\right) \ge (1-\delta) \nu (I(\mathcal{G}')).
    \end{equation*}
    We construct $\mathcal{T}$ as the union of three distinct trees $\mathcal T_1, \mathcal T_i$, and $\mathcal T_\infty$. In the first step we construct $\mathcal{T}_1 \subset \mathcal{G} \cap \mathcal E$. Each $\mathcal E$-equivalence class $K$ such that $K \cap O(\mathcal G') \neq \emptyset$ contains a connected component of $\mathcal T_1$. This connected component is any spanning tree of $K \cap I(\mathcal G')$, which is connected by hypothesis, attached to a unique vertex in $K\cap O (\mathcal{G}')$ by a single vertex. The graph $\mathcal T_1$ already satisfies the conclusion ii). In the next step of the proof we extend its domain so that i) is met. 
    
    Let $A\subset I$ be the Borel subset consisting of those $x\in I$ such that $[x]_E \cap O (\mathcal{G}') = \emptyset$. Let $A = A_f \sqcup A_\infty$ be a Borel partition into unions of connected components of $\mathcal{G}[A]$ such that $\mathcal{G}[A_f]$ has finite connected components and $\mathcal{G}[A_\infty]$ has infinite connected components. 
    
    We let $\mathcal{T}_f \subset \mathcal G$ be defined as follows. For each connected component $K$ of $\mathcal G [A_f]$, the forest $\mathcal T_f$ contains a connected component which is given by any spanning subtree of $\mathcal G [K]$, plus one single edge attaching said subtree to $X \backslash A$. 
    
    To conclude the proof, our goal is to partition $A_\infty$ into finite $\mathcal G$-connected subgraphs adjacent to $X\backslash A$.  To this end, fix a Borel linear order\footnote{This may be substituted by a factor of iid argument using a tie-breaking induced by Uniform$[0,1]$ labels which are all distinct almost surely.} $\leq$ on $X$ and for each $x \in \partial_v^{\mathrm{out}} (A_\infty)$ define 
    \begin{equation}\label{eq:voronoi}
        V_{x} := \{ y \in A_\infty : \forall z \in \partial_v^{\mathrm{out}} (A_\infty), d_{\mathcal{G}} (x,y) \leq d _{\mathcal{G}} (z,y)\ \text{and in case of equality}\ x \leq z\}.
    \end{equation}
    The Borel family of sets $(V_{x})_{x\in \partial_v^{\mathrm{out}} (A_\infty)}$ can be understood as a Voronoi tessellation of $A_\infty$ with centers in the boundary $\partial_v^{\mathrm{out}} (A_\infty)$. If $y \in A_\infty$ we let $c(y) \in B$ be the unique element of $B$ such that $y\in V_{c(y)}$. If the $V_{x}$ were finite for a.e.\ $x\in X$, then the proof would be over. This would be the case, for instance, in the pmp setting. However, this is not guaranteed in the mcp setting. To conclude the proof we need the following claim. It shows that the $V_x$ can be locally edited so that only an arbitrarily small proportion of vertices is contained in an infinite $V_x$. 

    \begin{claim}\label{claim:partition}
        Suppose that $(V_x)_{x\in \partial_v^{\mathrm{out}} (A_\infty)}$ is a partition of $A_\infty \cup \partial_v^{\mathrm{out}} (A_\infty)$ such that, for each $x\in \partial_v^{\mathrm{out}} (A_\infty)$, the induced subgraph $\mathcal G[V_x]$ is connected and $x\in V_x$. Then, for any $\eta > 0$ there exists a partition $(W_x)_{x\in B}$ of $A_\infty \cup \partial_v^{\mathrm{out}} (A_\infty)$ such that, for each $x\in \partial_v^{\mathrm{out}} (A_\infty)$, the induced subgraph $\mathcal G[W_{x}]$ is connected, satisfies $x\in W_x$, and 
        $$
            \nu\left(y \in A_\infty :  |W_{c(y)}|=\infty\right)\leq \eta \nu (A_\infty).
        $$
    \end{claim}

    \begin{proof}[Proof of claim]
     For each $x\in \partial_v^{\mathrm{out}} (A_\infty)$ and $k\in \mathbb N$, let $$
     V_x^{\leq k}:= \{y\in V_x : d_{\mathcal G} (y,x) \leq k\},
     $$
     and denote its complement $V_x \backslash V_{x}^{\leq k}$ by $V_x^{> k}$. Note that by local-finiteness $\mathcal G [V_x^{>k}]$ has at most finitely many connected components. For simplicity of the exposition we will assume that each $\mathcal G[V_x^{>k}]$ is actually connected. The following proof can be easily adapted to the full generality treating each connected component in the same fashion. 
     
     Fix $k\in \mathbb N$ large enough such that $$
    \nu \left(\bigcup_{x\in \partial_v^{\mathrm{out}} (A_\infty)} V_x^{\leq k}\right) \geq (1 - \eta/3) \nu (A_\infty).
    $$
    Let us call the union in the above expression $V^{\leq k}$ and its complement $V^{>k} := A_\infty \backslash V^{\leq k}$.
    
    Let $B \subseteq \partial_v^{\mathrm{out}} (A_\infty)$ be the Borel subset of those $x\in \partial_v^{\mathrm{out}} (A_\infty)$ for which $|\partial_v^{\mathrm{out}} (V_x^{>k}) \cap V^{\leq k} | = \infty$. By one-endedness and local finiteness, if $x \in B$ then $\partial_v^{\mathrm{out}} (V_x^{>k}) \cap V^{\leq k}$ intersects infinitely many other $V_y$ non-trivially. It follows that there exists $C \subseteq \partial_v^{\mathrm{out}} (A_\infty)$ such that every $x \in B$ admits a $y \in C$ such that $\partial_v^{\mathrm{out}} (V_x^{>k}) \cap V_y^{\leq k} \neq \emptyset$ and $C$ has small measure:
    $$
    \nu \left(\bigcup_{x\in C} V_x\right) \leq \frac{\eta}{3} \nu (A_\infty).
    $$
    Using the Lusin--Novikov Theorem, construct a measurable choice $y(x) \in C$ for each $x\in B$ such that $\partial_v^{\mathrm{out}} (V_x^{>k}) \cap V_{y(x)}^{\leq k} \neq \emptyset$. For each $x\in B$ we let the set $W_x$ to be $V_x^{\leq k}$ and place the remaining part $V_x^{> k}$ into $W_{y(x)}$. 
    
    % and such that every $x\in C$ has an $\mathcal H$-neighbour in $D$. For each $x \in C$, let $W_x := V_x^{\leq k}$, and for each $x \in D$ let $W_x$ be obtained by attaching the $V_y^{> k}$ of those $y$ such that $s (y) = x$, where $s$ is a Lusin--Novikov uniformization of $\mathcal H \cap (D \times C)$.

    Let $D := \partial_v^{\mathrm{out}} (A_\infty) \backslash B$. We define a Borel graph $\mathcal H\subseteq\mathcal{R}_{\mathcal{G}}$ on $D$ by letting $(y,x) \in \mathcal H$ if and only if $\partial_v^{\mathrm{out}} (V_x^{> k}) \cap V^{> k}_y \neq \emptyset$. By one-endedness of $\mathcal{G}$ the graph $\mathcal{H}$ is aperiodic. Since $\mathcal G$ is hyperfinite, so is $\mathcal H$. Thus, there exists a Borel subgraph $\mathcal K \subset \mathcal H$ with finite connected components and a Borel transversal $D' \subset D$ such that 
    $$
    \nu \left(\bigcup_{x\in D'} V_x\right) \leq \frac{\eta}{3} \nu (A_\infty).
    $$
    For each $x\in D$ we let $z(x) \in D'$ be the unique element of $D'$ which is $\mathcal K$-connected to $x$. For each $x\in D \backslash D'$ we let $W_x$ be $V_x^{\leq k}$ and place the remainder $V_x^{>k}$ in $W_{z(x)}$ where $y \in D'$ the unique element of the transversal in the $\mathcal K$-connected component of $x$. 

    The above process concludes with a partition $(W_x)_{x\in \partial_v^{\mathrm{out}} (A_\infty)}$ that satisfies the desired properties.

    
    \end{proof}

   To conclude the proof of the lemma set $V_x^{(1)} := V_x$ for each $x\in \partial_v^{\mathrm{out}} (A_\infty)$, and for each $n\in \mathbb N$ given $(V_x^{n})_{x \in \partial_v^{\mathrm{out}} (A_\infty)}$, construct $(V_x^{n+1})_{x \in \partial_v^{\mathrm{out}} (A_\infty)}$ applying the claim with $\eta = 2^{-n}$. By the Borel--Cantelli lemma, for $\nu$-a.e.\ $x\in (V_x^{n})_{x \in \partial_v^{\mathrm{out}} (A_\infty)}$ there exists a least $n_x \in \mathbb N$ such that $V_x^{(N)} = V_x^{(n_x)}$ for every $N \geq n_x$. The partition $(V_x^{(n_x)})_{x \in \partial_v^{\mathrm{out}} (A_\infty)}$ satisfies the desired properties.  
\end{proof}

\end{document}
